\documentclass[10pt]{article}
\usepackage[utf8]{inputenc}
% Language and font encodings
\usepackage[french]{babel}
\usepackage{concmath}
\usepackage[OT1]{fontenc}

% Sets page size and margins
\usepackage[letterpaper, margin=1in]{geometry}
% Useful packages
\usepackage{amsmath}
\usepackage{graphicx}
\usepackage{subcaption}
\usepackage{float}
\usepackage{booktabs}
\usepackage{svg}
\usepackage{multirow}
\usepackage{multicol}
\usepackage{enumitem}
\usepackage{listings}
\usepackage{color}
\usepackage{tikz}
\usepackage{pgfplots}
\usepackage{pgfplotstable}
\usepackage{caption}
\usepackage{wrapfig}
\usepackage{lipsum}
\usepackage{fancyhdr}
\usepackage{wrapfig}
\usepackage{indentfirst}
\usepackage{hyperref}
\usepackage{csquotes}
\usepackage{amsfonts}
\usepackage{tikz}
\usepackage{pgffor}
\usetikzlibrary{patterns}
\usetikzlibrary{intersections}
\usepackage{relsize}
\usepackage{extpfeil}
\usepackage{mathtools}
\usepackage{pgfplots}
\pgfplotsset{compat=1.18}
\usepackage{amsmath}
\usepackage{tikz}
\usetikzlibrary{arrows.meta}
%indentfirst package is used to indent the first paragraph of each section
\usepackage{indentfirst}

% Line spacing for clear readability
\usepackage{setspace}
\setstretch{1.25}

\usepackage{mathtools}
\DeclarePairedDelimiter\abs{\lvert}{\rvert}%
\DeclarePairedDelimiter\norm{\lVert}{\rVert}%
\DeclarePairedDelimiter\ceil{\lceil}{\rceil}
\DeclarePairedDelimiter\floor{\lfloor}{\rfloor}

% Custom environments (in French)
\usepackage{amsthm}            % For custom theorem-like environments
\newtheoremstyle{frenchstyle}  % Custom theorem style with no indent and new line for title
    {10pt}                     % Space above
    {10pt}                     % Space below
    {}                         % Body font (empty for normal font)
    {}                         % Indent (no indent)
    {\bfseries}                % Theorem head font
    {}                         % Punctuation after theorem head
    {\newline}                 % Space after theorem head
    {}                         % Theorem head spec (can be left empty)

\theoremstyle{frenchstyle}
\newtheorem{question}{Question}
\newtheorem{solution}{Solution}
\newtheorem{definition}{Définition}
\newtheorem{exercice}{Exercice}
% Idea developpement newtheorem
\newtheorem{idee}{Idée}

\newcommand{\R}{\mathbb{R}}
\newcommand{\rel}{\mathcal{R}}
\newcommand{\Z}{\mathbb{Z}}
\newcommand{\N}{\mathbb{N}}
\newcommand{\B}{\mathcal{B}}
% Paragraph settings

\setlength{\parskip}{1em}
\renewcommand{\thefootnote}{\roman{footnote}}
% indentation priority

\begin{document}

\begin{titlepage}
	\vspace*{9em}{\centering\Huge\usefont{T1}{qzc}{m}{it}
		DM MAT 328\par}
	\vspace*{1em}{\centering\Large\usefont{T1}{qzc}{m}{it} Devoir à la Maison\par}
	\vspace*{\fill}
	% Left Bottom Author informations
	\begin{minipage}[b]{.4\textwidth}
		\raggedright
		\scriptsize
		\textit{OZEL, Furkan}\\
		\textit{20504325}\\
		\textit{Galatasaray University, 6iéme Sem.}\\
		\textit{MAT 328}\\
		\
	\end{minipage}
\end{titlepage}

\tableofcontents
\newpage

\begin{exercice}
	\addcontentsline{toc}{section}{Exercice 1}
	D'abord tracer approximativement les lignes caractéristiques de l'équation du transport inhomogène suivante :
	\[
		u_t + \cos^2\left(\frac{x\pi}{2}\right)u_x = 0
	\]
	On donne les deux types de conditions initiales suivantes :
	\[u(0,x) = 1-e^{-\frac{1}{x^2}} \]
	A partir de la connaissance des lignes caractéristiques, donner le profil approximatif des solutions de l’équation ci-dessus à des temps $t > 0$ différents.

	Nous allons maintenant résoudre explicitement :
	\begin{enumerate}
		\item Donner grâce à la fonction $\beta$ l’expression exacte de ces lignes caractéristiques, en fonction de la constante d’intégration $\kappa$ (kappa) ou l’ordonnée à l’origine $k$;
		\item Résoudre ensuite explicitement le problème grâce aux deux types de conditions initiales suivantes :
		      \[\text{1. } u(0, x) = 1-e^{-\frac{1}{x^2}} \quad \text{puis} \quad \text{2. } u(1, x) = 1-e^{-\frac{1}{x^2}} \]
	\end{enumerate}

	Maintenant on étudie l’équation du transport inhomogène avec amortissement suivante :

	\[u_t + \cos^2\left(\frac{x\pi}{2}\right)u_x + u=0\]

	Résoudre avec les conditions initiales $u(0,x) = 1-e^{-\frac{1}{x^2}}$.
\end{exercice}

\begin{solution}
	\addcontentsline{toc}{subsection}{Solution de l'exercice 1}
	On a une equation du transport inhomogène de la forme suivante:
	\[
		u_t + \cos^2\left(\frac{x\pi}{2}\right)u_x = 0
	\]

	C'est possible que voir le vitesse d'onde $c(x) = \cos^2\left(\frac{x\pi}{2}\right)$ est une fonction périodique. Par consequent,
	\[
		\frac{dx}{dt} = c(x) = \cos^2\left(\frac{x\pi}{2}\right)
	\]

	Examinons la fonciton $c(x) = \cos^2\left(\frac{x\pi}{2}\right)$, c'est une fonciton trigonométrique qui est périodique avec une périodique.  On sait que $c(x) \in [0,1]$, on les cas différents les suivantes :
	\begin{itemize}
		\item $c(x) = 0$ pour $x \in 2\Z$.
		\item $c(x) = 1$ pour $x \in 2\Z + 1$.
		\item $c(x) \in (0,1)$ pour $x \in \R \setminus \left( 2\Z \cup 2\Z + 1 \right)$.
	\end{itemize}

	On peut essayer l'equation différentielle ordinaire suivante:
	\[
		\frac{dx}{dt} = \cos^2\left(\frac{x\pi}{2}\right)
	\]

	Alors,
	\begin{align*}
		\int \frac{dx}{\cos^2\left(\frac{x\pi}{2}\right)} & = \int dt                                      \\
		\int \frac{2}{\pi \cos^2(u)} du                   & = \int dt \quad \textit{; $u=\frac{\pi x}{2}$} \\
		\frac{2}{\pi} \int \sec^2(u) du                   & = t + \kappa                                   \\
		\frac{2}{\pi} \tan(u)                             & = t + \kappa                                   \\
		\frac{2}{\pi} \tan\left(\frac{x\pi}{2}\right)     & = t + \kappa                                   \\
	\end{align*}

	$\kappa$ est la constante d'intégration qui est la valeur de $x$ à $t=0$. Par consequent,
	\[
		\tan(\frac{x\pi}{2}) = \frac{\pi}{2}(t + \kappa) \quad \text{ou} \quad x = \frac{2}{\pi}\arctan\left(\frac{\pi}{2}(t + \kappa)\right)
	\]

	On a obtenue la formula $x(t)$ des lignes caractéristiques. Supposons que $x_0$ est la valeur de $x$ à $t=0$, alors on a $x(0) = x_0$. Par consequent, on a $\kappa = \frac{2}{\pi}\tan\left(\frac{x_0\pi}{2}\right)$.

	De plus, on sait que $\frac{du}{dt}=0$. Cela signifie que la solution de l'équation du transport inhomogène est constante le long des lignes caractéristiques. Par consequent, on a $u(x(t),t)=u(x_0,0)$.

	On peut essayer de dessiner les lignes caractéristiques. D'abord examinons les propriétés des lignes caractéristiques. On a trouvé que les lignes caractéristiques sont données par la formule suivante:
	\[
		x(t) = \frac{2}{\pi}\arctan\left(\frac{\pi}{2}(t + \kappa)\right).
	\]

	\begin{figure}[h]
		\centering
		\includesvg[width=0.5\textwidth]{karakteristik_baslangidddc_renkli.svg}
		\caption{Lignes caractéristiques de l'équation du transport inhomogène.}
		\label{fig:lines_characteristics}
	\end{figure}

	On sait que la solution \( u \) est constante le long des lignes caractéristiques. En particulier, si on note \( x_0 \) la position initiale telle que \( x(0) = x_0 \), alors \( u(x(t), t) = u(x_0, 0) \).

	Puisque nous connaissons la relation entre \( x \), \( t \) et \( x_0 \), on peut écrire :
	\[
		x = \frac{2}{\pi} \arctan\left( \frac{\pi}{2}(t + \kappa) \right) \Rightarrow \kappa = \frac{2}{\pi} \tan\left( \frac{x\pi}{2} \right) - t.
	\]
	On déduit alors la position initiale \( x_0 \) à partir de \( \kappa \) :
	\[
		x_0 = \frac{2}{\pi} \arctan\left( \frac{\pi}{2} \kappa \right) = \frac{2}{\pi} \arctan\left( \frac{\pi}{2} \left( \frac{2}{\pi} \tan\left( \frac{x\pi}{2} \right) - t \right) \right).
	\]

	Donc la solution explicite avec la condition initiale \( u(0,x) = 1 - e^{-1/x^2} \) est donnée par :
	\[
		u(x, t) = 1 - \exp\left( -\frac{1}{\left[ \frac{2}{\pi} \arctan\left( \frac{\pi}{2} \left( \frac{2}{\pi} \tan\left( \frac{x\pi}{2} \right) - t \right) \right) \right]^2} \right).
	\]

	Ainsi, on a obtenu la solution explicite du problème de Cauchy pour l'équation du transport inhomogène avec condition initiale donnée.

	Passons maintenant à la deuxième condition initiale, \( u(1, x) = 1 - e^{-1/x^2} \).

	Dans ce cas, la solution est toujours constante sur les lignes caractéristiques, donc :

	\[
		u(x(t), t) = u(x_0, 1) \quad \text{où} \quad x(1) = x_0
	\]

	La relation caractéristique devient :

	\[
		x = \frac{2}{\pi} \arctan\left( \frac{\pi}{2}(t + \kappa) \right), \quad \text{et} \quad x_0 = \frac{2}{\pi} \arctan\left( \frac{\pi}{2}(1 + \kappa) \right).
	\]

	D'où :

	\[
		\kappa = \frac{2}{\pi} \tan\left( \frac{x\pi}{2} \right) - t \quad \Rightarrow \quad x_0 = \frac{2}{\pi} \arctan\left( \frac{\pi}{2} \left( 1 + \frac{2}{\pi} \tan\left( \frac{x\pi}{2} \right) - t \right) \right)
	\]

	et donc :

	\[
		u(x, t) = 1 - \exp\left( -\frac{1}{\left[ \frac{2}{\pi} \arctan\left( \frac{\pi}{2} \left( 1 + \frac{2}{\pi} \tan\left( \frac{x\pi}{2} \right) - t \right) \right) \right]^2} \right)
	\]

	Cela termine la résolution explicite de l'équation du transport inhomogène pour les deux types de conditions initiales proposées.
\end{solution}

\newpage

\begin{exercice}
    \addcontentsline{toc}{section}{Exercice 2}
	On considère l’équation du transport
	\begin{equation}
		u_t + c(t, x)u_x = 0
	\end{equation}

	où la vitesse de l'onde $c(t ,x)$ dépend maintenant des deux variables temps est espace.

	\begin{enumerate}
		\item On définit les courbes caractéristiques comme toutes les courbes paramétrées du plan $\R^2$ de coordonnées $(t, x)$ qui satisfont l'équation différentielle ordinaire :
		      \[
			      \frac{dx}{dt} = c(t, x).
		      \]
		      Montrer que toute solution $u(t ,x)$ de l'équation $(1)$ est constante sur les lignes caractéristiques.
		\item Déterminer les lignes caractéristiques de l’équation (on rappelle que chaque ligne caractéristique est étiquetée par une constante d’intégration) :
		      \begin{equation}
			      u_t + t^2u_x = 0
		      \end{equation}
		\item Montrer que la fonction $\xi(t,x) = x - \frac{t^3}{3}$ est constante le long des lignes caractéristiques. \textit{on peut dériver $\xi$ par le vecteur $\vec{L}$ tangent aux lignes caractéristiques.}
		\item On pose $v(t, \xi) = u(t, x)$. Montrer que la fonction $v$ ne  dépend que de la variable $\xi = x - \frac{t^3}{3}$. Par la suite on note $v(t, \xi) = v_0(\xi)$ pour une certaine fonciton $v_0$, qui sera fixée par les conditions initiales.
		\item On prend comme condition initiale $u(0,x)=e^{-\frac{1}{x^2}}$. Résoudre l'équation différentielle $(1)$.
	\end{enumerate}
\end{exercice}

\begin{solution}
	\addcontentsline{toc}{subsection}{Solution de l'exercice 2}
    L'équation de transport, qui est une EDP quasi-linéaire du premier ordre, aura la forme suivante :
    \begin{equation*}
        a(t, x, u)u_t + b(t, x, u)u_x = f(t, x, u)
    \end{equation*}

    mais on a une équation avec $a(t, x, u) = 1$, $b(t, x, u) = c(t, x)$ et $f(t, x, u) = 0$. Ce satisfait la condition d'être une EDP quasi-linéaire du premier ordre.
	
	On veut montrer que $u(t, x)$ est constant le long des lignes caractéristiques,
    \begin{align*}
        \frac{dx}{dt} = c(t,x)
    \end{align*}

    Supposons que $u(t, x)$ est une solution de l'équation de transport. Prendrons une ligne characteristique et écrirons $(t,x(t))$ liéu à $t$ parametre, par consequent $\exists x(t), \frac{dx}{dt}=c(t,x(t))$. On a la relation suivante:
    \[
        U(t)=u(t,x(t))
    \]

    Calculons les derivatives de $U(t)$ en utilisant la règle de la chaine:
    \begin{align*}
        \frac{dU}{dt}&=\frac{\partial u}{\partial t}(t,x(t))\frac{dt}{dt}+\frac{\partial u}{\partial x}(t,x(t))\frac{dx}{dt}\\
        &\implies \frac{dU}{dt} = u_t(t,x(t)) + u_x(t,x(t))c(t,x(t))\\
    \end{align*}

    Puisque $u(t,x)$ est une solution de l'équation de transport $u_t+c(t,x)u_x=0$, cette équation est satisfaite pour tous les points $(t,x)$ dans le domaine de la solution. En particulier, elle est satisfaite pour les points $(t,x(t))$ qui se trouvent sur une courbe caractéristique. $u+(t,x(t))+c(t,x(t))u_x(t,x(t))=0$On peut maintenant substituer cette relation dans l'expression de $\frac{dt}{dU}$ obtenue à l'étape précédente :
    \begin{align*}
        \frac{dU}{dt} &= u_t(t,x(t)) + u_x(t,x(t))c(t,x(t))\\
        &= 0
    \end{align*}
    
    Cela signifie que la fonction $U(t)$ est constante le long de la courbe caractéristique. Par conséquent, $u(t,x)$ est constant le long des lignes caractéristiques.

    On a $u_t + t^2u_x = 0$ avec $c(t,x)=t^2$, cette fois n'est pas liéu aux espaces. On peut definir les lignes characteristique en utilisant ODE : 
    \[
        \frac{dx}{dt}=c(t,x)
    \]

    Alors,
    \begin{align*}
        \frac{dx}{dt}&=c(t,x)\\
        \int dx = t^2dt\\
        x(t) &= \frac{t^3}{3} + \kappa\\
    \end{align*}

    Par consequent,
    \[
        \kappa = x - \frac{t^3}{3}
    \]

	Les lignes caractéristiques oeuvent étre dessinées comme suit :
	\begin{figure}[h]
		\centering
		\includesvg[width=0.6\textwidth]{deues.svg}
	\end{figure}

    On veut montrer que $\xi(t,x)=x-\frac{t^3}{3}$ est constante sur les lignes caracteristique de $u_t + t^2u_x=0$, on peut trouver le vecteur qui est tengente des lignes caractéristiques.

    On a developpé a idee pour $a(t,x)u_t + b(t,x)u_x=0$ à liéu $L=(a(t,x),b(t,x))$ qui est area du vecteur caractéristique. Nos equations satisfait que les suivantes : 
    \begin{align*}
        a=1, b=1 &\implies L = (1, c(t,x))\\
        &\implies L=(1, t^2)
    \end{align*}

    Pour qu'une fonction soit constante le long des courbes caractéristiques, cela signifie exactement que sa dérivée dans la direction du champ de vecteurs caractéristique $L$ est nulle. Alors, en utilisant le formule suivante:
    \[
        L[F]=\frac{\partial F}{\partial t} + t^2 \frac{\partial F}{\partial x}
    \]

    Calculons,
    \begin{align*}
        \frac{\partial \xi}{\partial t} &= \frac{\partial}{\partial t}\left(x - \frac{t^3}{3}\right) = 0 - \frac{3t^2}{3} = -t^2\\
        \frac{\partial \xi}{\partial x} &= \frac{\partial}{\partial x}\left(x - \frac{t^3}{3}\right)=1-0=1
    \end{align*}

    Raplaçons dans la formule de $L$:
    \begin{align*}
        L[\xi] &= \frac{\partial \xi}{\partial t} + t^2 \frac{\partial \xi}{\partial x}\\
        &= -t^2 + t^2(1)\\
        &= 0
    \end{align*}

    Alors, $L[\xi] = 0$ cela signifie que $u(t,x)$ est constant le long des lignes caractéristiques. 

    En fait on a une system comme suivante :
    \begin{align*}
        \begin{cases*}
            u_t + c(t,x)u_x=0 \quad x\in \R, r>0\\
            u(0,x)=\exp(-\frac{x^2}{2})
        \end{cases*}
    \end{align*}

    On pose maintenant $v(t, \xi) = u(t, x)$ avec $\xi = x - \frac{t^3}{3}$. Cela signifie que $u(t, x)$ peut être réécrit en termes de $v$ comme suit :
    \[
    u(t, x) = v(t, \xi) = v\left(t, x - \frac{t^3}{3}\right).
    \]

    Calculons les dérivées partielles de $u$ en fonction de $t$ et $x$ :
    \[
    u_t = \frac{\partial v}{\partial t} + \frac{\partial v}{\partial \xi} \frac{\partial \xi}{\partial t}, \quad u_x = \frac{\partial v}{\partial \xi} \frac{\partial \xi}{\partial x}.
    \]

    Or, on sait que $\xi = x - \frac{t^3}{3}$, donc :
    \[
    \frac{\partial \xi}{\partial t} = -t^2, \quad \frac{\partial \xi}{\partial x} = 1.
    \]

    Ainsi :
    \[
    u_t = \frac{\partial v}{\partial t} - t^2 \frac{\partial v}{\partial \xi}, \quad u_x = \frac{\partial v}{\partial \xi}.
    \]

    Substituons ces expressions dans l'équation de transport $u_t + t^2 u_x = 0$ :
    \[
    \left(\frac{\partial v}{\partial t} - t^2 \frac{\partial v}{\partial \xi}\right) + t^2 \frac{\partial v}{\partial \xi} = 0.
    \]

    Les termes en $\frac{\partial v}{\partial \xi}$ s'annulent, ce qui donne :
    \[
    \frac{\partial v}{\partial t} = 0.
    \]

    Cela montre que $v$ ne dépend pas de $t$, donc $v(t, \xi) = v_0(\xi)$, où $v_0$ est une fonction de $\xi$ seule. Par conséquent, la solution de l'équation de transport est donnée par :
    \[
    u(t, x) = v_0\left(x - \frac{t^3}{3}\right).
    \]

    Pour déterminer $v_0$, utilisons la condition initiale $u(0, x) = e^{-\frac{1}{x^2}}$. À $t = 0$, on a $\xi = x$, donc :
    \[
    v_0(x) = e^{-\frac{1}{x^2}}.
    \]

    Ainsi, la solution finale est :
    \[
    u(t, x) = e^{-\frac{1}{\left(x - \frac{t^3}{3}\right)^2}}.
    \]



\end{solution}

\end{document}